\honest{Is That Your True Rejection?}{Eliezer Yudkowsky}{2008}

Some people hear my ideas about superintelligence and Friendly AI and reject them. Asked why, they sometimes say, ``Why should I believe anything Yudkowsky says? He doesn't have a PhD!'' Some advise me to get one. There are good and bad reasons to get a PhD, and this is one of the bad ones.

My explanation of my critics is this. The idea reminds them of a category like ``strange weird idea,'' science fiction, an end-of-the-world cult or overenthusiastic youth, and they reject it ``at the speed of perception.'' Afterwards they look for a respectable reason and land on the PhD. From two pieces of evidence, both below, I describe what happens in their heads.

The ``true reason'' is whatever caused the rejection at the very first moment. The search for a reason ``won't necessarily hit on the true reason.'' I still ask people to report it, and give no way to tell when they have.

Nobody asks for my PhD when I write about human rationality. \nb{But a reader can check claims about reasoning against their own experience, and cannot check claims about superintelligence. A year earlier, in ``Argument Screens Off Authority,'' Yudkowsky had written that when a listener cannot evaluate an argument, the speaker's authority is what they have to go on.}

If I had a PhD, they would ask for tenure at a major university instead, and people do not believe Harvard professors who say weird things either. For a claim that sounds wrong to a novice, told by a stranger to someone outside the field, I suspect that credentials start to outweigh a first impression only at about the level of a Nobel prize. ``If that.''

My one example is Eric Drexler. People asked him for technical details, or told him to come back when he had a PhD. He spent six years writing the details, got his PhD under Marvin Minsky for them, and published them as \textsc{Nanosystems}, ``a great book.'' Did those people change their minds? ``Not so far as I ever heard.'' \nb{After the doctorate, Drexler's best-known critic was Richard Smalley, a Nobel laureate in chemistry, who argued in print that Drexler's machines could not work. The post does not mention him, or consider that the doubters stayed doubters for technical reasons.}

The idea is useful in business. When a venture capitalist says ``If only your sales were growing a little faster!'' or a customer says you lack feature X, that may not be the true rejection, and fixing it may change nothing. Think about that before spending great effort.

In disagreements, Robin Hanson and I believe that two rationalists should not agree to disagree: they should not have common knowledge of a disagreement unless something is very wrong. If a disagreement survives the first exchange, its true sources are probably hard to communicate or hard to expose. If they could all be laid on the table easily, it would have been settled at the first meeting. I list them: rare but well-supported knowledge or math, long inferential distances, hard-to-state intuitions, the outlook of a profession, patterns recognized from experience, habits of thought, emotional commitment to an outcome, fear that a past mistake could be exposed, and self-deception for the sake of pride. I give no way to tell which one applies in a given case.

My advice: ask yourself whether this is your true rejection. Openly psychoanalyzing the other person makes a conversation go bad ``very fast.'' I have spent most of this post doing that to my critics. In the last paragraph I decide that asking the other person is ``fair game'' after all, if done humbly, as long as the most embarrassing possibilities are left unspoken. The question I propose is: ``Is that simple straightforward-sounding reason your true rejection, or does it come from intuition-X or professional-zeitgeist-Y?'' I have just listed them.

\nb{Two days later Yudkowsky asked Robin Hanson, in public and ``in all humility,'' whether this was his true rejection, as the last paragraph allows.}
